\section{Validity and limits of generalization}
\label{sec:heldout-limits}
The initial and expanded factorial studies use two requested model aliases, each at one reasoning setting and one client version. Their samples and replication counts differ. This provides more scope than a single model run but does not identify model effects or generalize across model families. The development repeats are a separate feasibility study. Public task availability leaves training contamination unresolved; related tasks can also violate an independent-task approximation. A different model, language, task distribution or service revision could reverse an observed contrast. The recorded client version and execution dates do not fix the provider's model weights.

The deadline and administrative continuation create a second boundary. Failed generations remain missing, and the amendment is disclosed. Inference from available pairs concerns the observed eligible subset. Assignment-denominator ranges show extreme missing-outcome possibilities but do not correct selection bias. A timeout could reflect computational difficulty rather than random transport noise, so failure to reject a quality difference does not prove equivalence. No externally justified quality-loss margin or joint quality--cost criterion supports a non-inferiority or economic-superiority decision.

The interventions concern explicit instructions. Role labels do not certify independent agents, and call boundaries also expose intermediate responses. Although task information and stage operations are controlled, there is no matched hard output-token allowance. Resource comparisons include induced reasoning length, repeated context transmission and client overhead. List-price valuation depends on cache state; the no-cache sensitivity removes the discount but does not make subscription execution identical to an API deployment. Research and evaluator compute are excluded and no total cost of ownership is estimated.

Native tests improve auditability while retaining benchmark limitations. Gold and incorrect controls detect an unusable environment or an insensitive negative check, but passing them does not prove that every original assertion implements the natural-language task. A separate development audit records instruction/test disagreements as annotations; the original prompts, tests, and scores remain unchanged. These annotations are not post-hoc exclusions or a replacement gold standard. The seven reference-ineligible tasks in the initial mini study and the 15 in the Luna study remain in their respective assignment and resource accounting.

The INoT comparison is a disclosed prompt-level adaptation with an explicit resolution of source ambiguity. Its separate batch retains timing and cache differences. The completed Luna extension is a separate primary study with 15,000 assigned generations, 14,961 complete candidates and 985 control-eligible tasks; it reuses the initial 200 task IDs and therefore is not a new held-out sample or fixed-weight replication. Self-Collaboration (SCC) is being compared in an amended 2,000-block prefix (6,000 assigned method cells); generation is ongoing and native outcomes have not been inspected. Eight interrupted submitted attempts remain paused unknowns and are not retried. Native author INoT, Agentless and SWE-agent implementations are also not evaluated. The latter systems introduce retrieval, tools and repository interaction outside this fixed-context factorial question. The one-issue SWE-bench demonstration provides genuine patch and test evidence, but its complete ten-assignment matrix on a single issue cannot support a repository-benchmark score or comparative agent ranking. A native multi-repository comparison and a matched-budget, multi-model replication remain necessary for those broader claims.
